\section{Conclusion}
%In this paper, we provide a comprehensive look at 
This paper presents the design, implementation, and deployment of \sys{}, a production-grade checkpointing system built for LFM development. 
\sys{} advocates a unified architecture for %both storage representation and 
checkpoint representation and workflow, enables efficient load-time  resharding and supports multiple training frameworks and storage backends.
It enables full-stack I/O performance and scalability optimizations and incorporates efficient performance monitoring and analysis tools for LFM development at scale.
Compared to state-of-the-art baselines~\cite{torch-dcp, megatron-dcp}, \sys{} %achieves significant performance improvement, 
reduces checkpoint stalls by up to 161.50$\times$ and shortens end-to-end checkpointing completion time by up to 9.96$\times$. Additionally, the loading (resharding) procedure is optimized to be 3.88$\times$ faster on average.
We believe our system not only provides valuable practical experience to those building checkpointing systems for real-world LFM development but also offers deep insights for future research in the community.
%We propose \sys{}, a PyTorch-native LLM checkpointing system designed for efficient %automatic load-time checkpoint resharding and flexible multi-framework support.
%\sys{} utilizes a disaggregated storage architecture for checkpoint data and metadata, %effectively decoupling the checkpoint storage from various parallelism configurations and %training frameworks.
%\sys{} addresses the challenges of irregular tensor resharding through the use of %asynchronous tensor merging and enhances the stability of planning by implementing a tree-%based communication topology.
%Additionally, we propose several I/O performance optimization techniques aimed at boosting %system efficiency.
%These techniques include fine-grained save/load pipelines, the Ping-Pong memory pool, %workload-balanced saving, zero redundancy loading, etc.
%Extensive Experimental results demonstrate the superiority of \sys{} over baselines on %various real-world checkpoint saving and loading tasks.

% \yhpeng{polish the references using unified format, pay attention to the case sensitivity in the bib title, and sort the references according to cited order}

% \yhpeng{To limit the paper to 12 pages, (1) move dataloader and extra states, dataloader state prefetching to Appendix; (2) move Correctness verification in the Evaluation to Appendix; (3) For Sec. 2.2.1 and Fig.3, move the discussion about offline scripts to Appendix; (4) move integrity check to Appendix; (5) Fig.7 can be deleted and just use Fig.8 or Fig.9; (6) For Fig.8, only drawing 2 ranks is ok, no need of 4 ranks, (7) Table 10 in the related work can be deleted.}